% Shared macros for SR_algorithms.tex and every file in sections/.
% Only these packages are available (TinyTeX): amsmath, amssymb, booktabs, longtable, enumitem,
% xcolor, hyperref, float, geometry, tikz, caption. NOT available: algorithm2e, algorithmicx, microtype.
% Pseudocode = numbered enumerate lists.

\definecolor{srccol}{RGB}{90,90,140}
\definecolor{opencol}{RGB}{190,30,30}

% Inline code / identifiers (underscores are safe inside).
% \path (url package, loaded by hyperref) allows line breaks inside long identifiers.
\makeatletter\g@addto@macro\UrlBreaks{\do\:\do\-\do\,}\makeatother
\Urlmuskip=0mu plus 1mu
\urlstyle{tt}
\newcommand{\code}[1]{\ifmmode\text{\path{#1}}\else\path{#1}\fi}
% Source citation for an equation or statement: \src{SRK/fp_likelihood.py:fpt_full}
\newcommand{\src}[1]{{\small\textcolor{srccol}{[\path{#1}]}}}
% Anything not verified against the source, or a question for the author. Collected in the appendix.
\newcommand{\openq}[1]{{\color{opencol}[\textbf{OPEN:} #1]}}

% Model notation (physical)
\newcommand{\xc}{x_c}
\newcommand{\hext}{h_{\mathrm{ext}}}
\newcommand{\eps}{\varepsilon}
% Dimensionless groups (as named in the code)
\newcommand{\reta}{\rho_\eta}
\newcommand{\rbeta}{\rho_\beta}
\newcommand{\reps}{\rho_\varepsilon}
% Probability / survival
\newcommand{\dd}{\mathrm{d}}
\newcommand{\E}{\mathbb{E}}

% Per-algorithm template headings (use in this order; omit none, write "None known." if empty)
\newcommand{\tplPurpose}{\paragraph{Purpose.}}
\newcommand{\tplMath}{\paragraph{Mathematics.}}
\newcommand{\tplImpl}{\paragraph{Discretisation and implementation.}}
\newcommand{\tplLimits}{\paragraph{Assumptions and validity limits.}}
\newcommand{\tplWhere}{\paragraph{Where it lives today.}}
\newcommand{\tplEvidence}{\paragraph{Validation evidence.}}
\newcommand{\tplIssues}{\paragraph{Known issues and planned changes.}}
\newcommand{\tplRefs}{\paragraph{References.}}

% A verify marker (spec 004 FR-020, constitution II): \verify{results JSON}{dotted path}{format spec}{text}. It prints
% the text; a checker recomputes format(value at the path, spec) from the stored result and compares it with the text,
% as tools/verify_readme.py does for the README's <!-- verify:path:key:spec -->text<!-- /verify --> (path elements:
% dict keys or [i] list indices; an empty spec is str(value)).
\newcommand{\verify}[4]{#4}
